\chapter{Анализ подходов к построению вопросно-ответных систем для корпоративной базы знаний}
\label{chapter:analytics}

\section{Анализ предметной области RobyMAX}

В рамках проекта RobyMAX рассматривается корпус русскоязычных бизнесово-технических документов, связанных с продуктами, техническими требованиями, пресейлом, внедрением, сервисной поддержкой и коммерческими условиями. В тексте работы RobyMAX рассматривается как предметная область корпоративной базы знаний; производственные внедрения, реальные клиенты и коммерческий эффект не утверждаются, поскольку такие данные не следуют из репозитория.

Особенность корпуса состоит в сочетании материалов разного типа. Продуктовые документы описывают решения и сценарии применения, технические документы фиксируют требования к помещению, инфраструктуре, сети и эксплуатации, пресейл-материалы связаны с квалификацией запроса и подготовкой внедрения, сервисные документы описывают поддержку и типовые неисправности, коммерческие материалы отражают порядок согласования условий. Пользовательский вопрос может относиться к любому из этих слоев, а ответ может находиться внутри раздела документа, таблицы или PDF-страницы.

По отчету \path{data/reports/parsing/parse_report.json} в реестре проекта зафиксировано 36 документов, все они включены в обработку и успешно распарсены. Корпус включает 18 файлов DOCX, 12 файлов PDF и 6 файлов XLSX. После парсинга получено 2300 структурированных блоков общим объемом 589074 символов. Это показывает, что даже учебный корпус уже является неоднородным по форматам, структуре и характеру представления информации.

\section{Проблема поиска информации в документах}

Поиск информации в такой базе знаний осложняется несколькими факторами. Во-первых, документы представлены в разных форматах, поэтому перед поиском требуется привести их к единому текстовому представлению. Во-вторых, пользователь обычно формулирует вопрос естественным языком, а не набором точных ключевых слов. В-третьих, релевантная информация может быть распределена между несколькими фрагментами. В-четвертых, пользователю требуется не только найти файл, но и получить прикладной ответ, связанный с источниками.

Классический пример такой проблемы — запрос о требованиях к инфраструктуре перед пилотным проектом. Ответ может затрагивать помещение, маршруты движения, сеть Wi-Fi, техническую приемку и фиксацию результатов пилота. Поиск только по названию файла или по одному ключевому слову не гарантирует, что будут найдены все существенные фрагменты. Поэтому единицей поиска целесообразно считать не документ целиком, а небольшой смысловой фрагмент, сохраняющий связь с исходным документом.

\section{Обзор существующих подходов}

Для решения задачи были рассмотрены четыре класса подходов: ручной поиск, классический полнотекстовый поиск, семантический поиск на основе embeddings и Retrieval-Augmented Generation. Сравнение приведено в таблице~\ref{tbl:approaches}.

\begin{table}[h]
\caption{Сравнение подходов к поиску и ответам по корпоративной базе знаний}
\label{tbl:approaches}
\centering
\small
\begin{tabular}{|p{0.18\linewidth}|p{0.25\linewidth}|p{0.28\linewidth}|p{0.20\linewidth}|}
\hline
Подход & Преимущества & Недостатки & Применимость для RobyMAX \\
\hline
Ручной поиск & Не требует разработки; эксперт сам интерпретирует документы & Медленный; зависит от знания структуры архива; плохо масштабируется & Допустим только для малого числа документов и разовых запросов \\
\hline
Полнотекстовый поиск & Быстрый; прост в реализации; эффективен при точном совпадении терминов & Слабо учитывает переформулировки; возвращает документы или фрагменты без готового ответа & Может быть полезен как дополнительный механизм, но не решает всю задачу MVP \\
\hline
Embedding search & Учитывает смысловую близость вопроса и фрагмента; работает с естественным языком & Требует выбора модели, построения индекса и оценки качества выдачи & Подходит как базовый способ поиска релевантных чанков в русскоязычном корпусе \\
\hline
RAG & Совмещает поиск и генерацию ответа; позволяет возвращать источники; база знаний обновляется без дообучения LLM & Качество ответа зависит от retrieval, чанкинга и prompt; требуется контролировать отсутствие сведений & Выбран как основной подход для первой MVP-версии \\
\hline
\end{tabular}
\end{table}

Полнотекстовые методы остаются важным классом информационного поиска~\cite{Manning2008IR,Robertson2009BM25}. Однако первая версия системы сознательно ограничивается dense retrieval, чтобы проверить базовый путь от вопроса к найденным фрагментам и ответу. Такое ограничение делает результаты первой итерации интерпретируемыми: можно отдельно оценить качество подготовки документов, embeddings, Qdrant-индекса и генерации ответа по найденному контексту.

\section{Обоснование выбора базового RAG}

RAG выбран как компромисс между поисковой системой и генеративной системой без внешнего контекста. Внешняя база знаний снижает риск неподтвержденных ответов, поскольку LLM получает конкретные фрагменты документов. Обновление базы знаний не требует дообучения модели: достаточно повторить обработку документов, построить embeddings и обновить индекс. Кроме того, система может возвращать \texttt{chunk\_id} и сведения об исходном документе, что делает ответ проверяемым.

Для первой MVP-версии важна завершенность минимального сценария. Базовый dense RAG включает кодирование вопроса embedding-моделью, поиск ближайших чанков в Qdrant, сбор top-k контекста, передачу вопроса и контекста в LLM, генерацию ответа на русском языке и вывод источников. Более сложные механизмы улучшения качества поиска не включаются в основной результат этой записки и рассматриваются как последующие этапы.

\section{Выводы по главе}

В результате анализа установлено, что задача поиска по базе знаний RobyMAX требует работы с разнородными документами, естественно-языковыми вопросами и фрагментами внутри документов. Ручной и классический полнотекстовый поиск не обеспечивают полноценного вопросно-ответного сценария. Для первой MVP-версии обоснован выбор базового RAG-подхода, включающего подготовку корпуса, dense retrieval в Qdrant и генерацию ответа по найденному контексту.

%%% Local Variables:
%%% TeX-engine: xetex
%%% End:
